\scanpage{69}
% archive sheet number: 68
% L1-3
$k$ corps, $k'$ extension finie

$S=\operatorname{Spec}k$, $S'=\operatorname{Spec}k'$

\margin{N.B. les \unsure{mêmes|énoncés} sont vrais : voir dans la \unsure{généralisation} ---}

\unsure{a)} $X'$ préschéma sur $k'$, $X_1=\prod_{S'/S}X'/S'$.

% L4
\[
\boxed{H^0(k,X_1)\simeq H^0(k',X)}\tag*{(i)}
\]
% L5-6
(ii) Si $X'$ est un schéma en groupes, soit $G'$, \uline{\unsure{lisse} sur $k$}, on a
\[
\boxed{H^1(k,G_1)\simeq H^1(k',G')}
\]
(Classes de fibrés principaux homogènes ---)

% L7-8
(iii) Si $G'$ est \deleted{\unsure{un}} un schéma en groupe \uline{comm.}, soit $\unsure{M'|M}$, on a
\[
\boxed{H^i(k,M_1)\simeq H^i(k',M')}\qquad\text{pour }\forall i
\]
(Cohomologie galoisienne\unsure{...|.})

% L9
(i) Est vrai par définition (\unsure{on s'en fout du corps},

% L10-18
(ii) Soit $P'$ fibré principal homogène sous $G'$, alors $P_1=\prod_{S'/S}P'$ est un schéma sur $k$, sur lequel $G'$ opère. On voit que $P_1$ est formellement fibré pr. hom. sous $G_1$, d'autre \deleted{\unsure{part il est}\ \illegible{2}} part utilisant la \unsure{lissité|simplicité} de $G'/k'$, on voit que $P_1\ne\varnothing$, \deleted{\unsure{que c'est un fibré principal}} \inserted{\unsure{donc est un fibré principal}} homogène.
% slanted
\margin{\unsure{Contre-exemple} si $G'$ est comm. \unsure{radiciel} ?? Si $G'$ est comm. \unsure{ou propre|non propre|non lisse}/\unsure{$k'$} ??}
Le foncteur $P'\mathrel{\unsure{\rightsquigarrow|\longmapsto}}P_1$ est \uline{pleinement fidèle}.
\margin{(i.e. $\unsure{H^1|H_1|H P_1}$ est \unsure{bijectif})}
Pour montrer que c'est une équivalence de catégories, un argument de descente ramène au cas où $P_1$ est trivial ---

% L19-22
(iii) Soit $T/S$ \deleted{fini plat} quelconque, alors on a
\[
C^*(T/S,M_1)\simeq C^*(T'/S',M')
\]
d'où
\[
H^{\unsure{\alpha|q}}(T/S,M_1)\simeq H^\alpha(T'/S',M'),
\]
d'où, si $T$ parcourt certains $S$-préschémas
\[
\varinjlim_T H^\alpha(T/S,M_1)\simeq\varinjlim_T H^\alpha(T'/S',M'),
\]
% L23-26
Or, si $T$ parcourt les rev. étales (\unsure{ou p.ex.} les \unsure{préschémas} étales séparés quasi-finis surjectifs sur $S$) les $T'/S'$ sont \uline{cofinaux} dans la \inserted{catégorie} \unsure{des revêtements|des recouvrements} rel. à $S'$.

D'où la conclusion

\margin{TSVP}
